\begin{center}\begin{tikzpicture}[font=\small,>=stealth,line join=round]
\draw[gray!40,rounded corners=3pt] (0,0) rectangle (14.2,8.3);\node at (7.1,7.700000000000001) {方法二：左右两数都变号，乘积却相同};
\begin{scope}[shift={(4.1,2.8)},x=1.5cm,y=1.15cm]\draw[->,gray] (-2,0)--(2,0) node[right]{$x$};\draw[->,gray] (0,-1)--(0,3) node[above]{$y$};\draw[blue,thick,solid,domain=-1.5:1.5,samples=180] plot(\x,{(\x)^2});\fill[orange!85!black] (-1,1) circle(2.2pt);\fill[orange!85!black] (1,1) circle(2.2pt);\end{scope}\node[anchor=west,align=left,text width=5.0cm] at (8.4,4.0) {图取 $f(x)=g(x)=x$\\[10pt]$x=1$：$1\cdot1=1$\\[10pt]$x=-1$：$(-1)(-1)=1$\\[10pt]$fg=x^2$ 为偶函数};
\node[text width=12.6cm,align=center] at (7.1,.65) {实际例子呈现两负号抵消；一般情况仍由任意 $a,b$ 的代数运算证明。};
\end{tikzpicture}\end{center}
